\originalchapter{序列比较与极值统计}
\originalsection{比对中的物理信息}
DNA字母顺序保存了历史信息，但相同片段也可能偶然出现。一个得分函数必须同时描述匹配收益、替换代价和插入删除代价。设两条序列为 $a_1\cdots a_n$、$b_1\cdots b_m$，替换得分 $s(a_i,b_j)$，每个空位罚分 $g>0$。全局比对得分 $F_{ij}$ 是两个前缀的最优得分，递推为
\[
F_{ij}=\max\{F_{i-1,j-1}+s(a_i,b_j),F_{i-1,j}-g,F_{i,j-1}-g\}.
\]
边界 $F_{i0}=-ig,F_{0j}=-jg$。最后一步只可能匹配两个末字母、删除第一个末字母或删除第二个末字母；删除该步后必须留下最优子问题，否则能以更优前缀改善当前结果。这一论证说明递推为何确实枚举全部可允许路径。

局部比对允许任何位置开始和结束，所以递推再与 $0$ 取最大，边界也置零。最终从全矩阵最大元回溯，而全局比对从 $F_{nm}$ 回溯。得分并列时必须保留并列前驱；只返回一个回溯路径不意味着比对唯一。空位开罚与延长罚不同的模型需要额外状态，本式不包含这种仿射罚分。

同源关系还需要统计显著性检验。评分越高不一定越可靠，因为搜索数据库越大，随机产生一次高分的机会也越大。将得分转成一个不考虑数据库大小的百分数，会丢失此处最重要的比较条件。

\originalsection{高斯最大值与Gumbel极限}
设 $r_i$ 独立，分布函数为 $F$，最大值 $M_N=\max_i r_i$。事件 $M_N\le x$ 等价于全部 $r_i\le x$，因而 $P(M_N\le x)=F(x)^N$。高分端用尾概率 $\overline F=1-F$ 表示，若 $N\overline F$ 有限，则 $(1-\overline F)^N\simeq\exp(-N\overline F)$。

对零均值、标准差 $\sigma$ 的高斯分布，分部积分给出 $\overline F(x)\sim\sigma\e^{-x^2/(2\sigma^2)}/(\sqrt{2\pi}x)$。选中心 $b_N$ 使 $N\overline F(b_N)=1$，尺度 $a_N=\sigma^2/b_N$。在 $x=b_N+a_Nz$ 附近，平方指数的一阶变化为 $z$，而尾部前因子相对变化趋零，因此
\[
P\!\left(\frac{M_N-b_N}{a_N}\le z\right)\longrightarrow\exp(-\e^{-z}).
\]
主尺度 $b_N\sim\sigma\sqrt{2\ln N}$，但 $\ln\ln N$ 修正影响有限数据拟合。最可能值、均值和分位点也相差 $O(a_N)$，不可在有限样本中一概等同。最小能量对应 $-M_N$，符号翻转后尾部方向必须一起翻转。

\originalsection{局部高分的尾部尺度}
无空位比对可看成随机得分增量的部分和。若平均增量为负，路径通常向下，却偶尔达到较高正值。令 $\lambda>0$ 满足 $\avg{\e^{\lambda s}}=1$，则 $\e^{\lambda S_n}$ 是均值保持的乘法过程。越过高度 $S$ 的概率因而具有 $\e^{-\lambda S}$ 的指数尺度，越界细节影响前因子。把约 $nm$ 个搜索起点合并，高分累计分布常近似为 $\exp[-Knm\e^{-\lambda S}]$，相应期望高分次数为 $Knm\e^{-\lambda S}$。

这个推导解释 $\lambda$ 和数据库长度的位置，但没有证明任意有空位、相关序列都服从同一公式。低复杂度序列、周期性疏水模式以及相关数据库条目会破坏独立背景。实际检索中的背景校正与数据库版本必须记录。
